\documentclass[11pt,a4paper]{article}
\usepackage{fontspec,amsmath,amssymb,amsthm,microtype}
\setmainfont{lmroman10-regular.otf}[BoldFont=lmroman10-bold.otf,ItalicFont=lmroman10-italic.otf,BoldItalicFont=lmroman10-bolditalic.otf]
\usepackage{xcolor}
\usepackage[most]{tcolorbox}
\usepackage{mathrsfs}
\providecommand{\lt}{<}
\providecommand{\gt}{>}
\setmonofont{DejaVuSansMono.ttf}[BoldFont=DejaVuSansMono-Bold.ttf,ItalicFont=DejaVuSansMono-Oblique.ttf,BoldItalicFont=DejaVuSansMono-BoldOblique.ttf,Scale=0.86]
\newfontfamily\leanfallback{STIXTwoMath-Regular.otf}[Scale=0.86]
\definecolor{leanrule}{gray}{0.62}
\definecolor{leanback}{gray}{0.962}
\definecolor{leantitle}{gray}{0.22}
\definecolor{leanmuted}{gray}{0.38}
\newtcolorbox{leanbox}[1]{enhanced,breakable,colback=leanback,colframe=leanrule,
  boxrule=0.4pt,arc=0.6mm,left=2.4mm,right=2.4mm,top=1.8mm,bottom=1.8mm,
  before skip=0.8em,after skip=1em,fonttitle=\small\bfseries,coltitle=leantitle,
  colbacktitle=leanback,titlerule=0.3pt,title={#1}}
\newlength{\leancharwidth}
\newenvironment{leancode}{\par\smallskip\begingroup\ttfamily\footnotesize\raggedright
  \settowidth{\leancharwidth}{M}\setlength{\parskip}{0pt}\setlength{\parindent}{0pt}
  \frenchspacing}{\par\endgroup\smallskip}
\newcommand{\leanline}[2]{\hangindent=\dimexpr#1\leancharwidth+4\leancharwidth\relax
  \hangafter=1\noindent\hspace*{#1\leancharwidth}#2\par}
\newcommand{\leansource}[1]{\par{\footnotesize\color{leanmuted}#1}\par}
\newcommand{\leanlabel}[1]{\par\smallskip\noindent{\small\bfseries #1}\par}
\usepackage[margin=25mm]{geometry}
\usepackage[unicode,hidelinks]{hyperref}
\hypersetup{pdfauthor={LangevinLean Project}}
\setlength{\emergencystretch}{3em}
\setlength{\parskip}{0.4em}
\setlength{\parindent}{0pt}
\allowdisplaybreaks[2]
\raggedbottom
\title{\texorpdfstring{Gibbs invariance of the physical-time particle phase law}{Gibbs invariance of the physical-time particle phase law}}
\author{LangevinLean Project}
\date{}
\begin{document}\maketitle
Draft mathematical exposition; these notes may contain mistakes.


The invariant density calculation is proved as an identity of phase laws: the complete physical Brownian phase process, started at its normalized Gaussian-times-Gibbs target, has that same target distribution at every nonnegative time.



The exact clock conjugacy transfers the full unit-temperature stationary phase law to the physical particle temperature. Position is unchanged, so the stationary parameter marginal remains the powered likelihood.


\section{\texorpdfstring{Transport the proved full phase stationarity through the exact clock}{Transport the proved full phase stationarity through the exact clock}}


The constructed selected flow for the sum potential preserves \(\mathcal N(0,I)\otimes\pi_Z\), by the shared kinetic stationarity theorem. Push this endpoint law through \(D_{1/\sqrt N}(v,z)=(v/\sqrt N,z)\). The product pushforward is exactly \(\mathcal N(0,N^{-1}I)\otimes\pi_Z\), namely the kinetic Gibbs target. The physical KIPLD flow was defined by this very time and velocity change. The almost-everywhere identity for scalar multiplication identifies its \(L^{2}\) phase representative with the pushed constructed representative. This proves invariance of the whole phase law, including velocity–position independence, rather than only the position marginal or a formal zero-generator calculation.



The unit-temperature kinetic flow has the full Gaussian-times-Gibbs stationary phase law. The proved clock-and-velocity conjugacy transports that law to thermal covariance \(I/N\) and the same position Gibbs marginal. Hence the physical KIPLD integral solution is stationary under its phase target.

\begin{leanbox}{Lean: \texttt{kipldBrownianFlow\_\allowbreak{}stationary}}
\leansource{\texttt{kipldBrownianFlow\_\allowbreak{}stationary} \quad KIPLMC/\allowbreak{}Stationarity.lean, lines 11–52}
\begin{leancode}
\leanline{0}{theorem\ kipldBrownianFlow\_stationary\ \{N\ dθ\ dx\ :\ ℕ\}\ (hN\ :\ 0<N)}
\leanline{4}{(U\ :\ IPLA.Joint\ dθ\ dx→ℝ)\ (m\ :\ ℝ)\ (hm\ :\ 0<m)\ (L\ γ\ :\ ℝ≥0)}
\leanline{4}{(hdiff\ :\ Differentiable\ ℝ\ U)\ (hsc\ :\ StrongConvexOn\ univ\ m\ U)}
\leanline{4}{(hLip\ :\ LipschitzWith\ L\ (gradient\ U))\ (T\ :\ ℝ)\ (hT\ :\ 0≤T)\ :}
\leanline{4}{letI\ :=\ ULA.isProbabilityMeasure\_gibbsMeasure\ (@IPLA.potential\ N\ dθ\ dx\ U)\ ((N:ℝ)*m)}
\leanline{6}{(mul\_pos\ (by\ exact\_mod\_cast\ hN)\ hm)\ (IPLA.potential\_differentiable\ hdiff)\ (IPLA.potential\_strongConvex\ hsc)}
\leanline{4}{let\ π\ :=\ ULA.gibbsMeasure\ (@IPLA.potential\ N\ dθ\ dx\ U)}
\leanline{4}{let\ hπ\ :=\ ULA.memLp\_id\_gibbsMeasure\ (@IPLA.potential\ N\ dθ\ dx\ U)\ ((N:ℝ)*m)}
\leanline{6}{(mul\_pos\ (by\ exact\_mod\_cast\ hN)\ hm)\ (IPLA.potential\_differentiable\ hdiff)\ (IPLA.potential\_strongConvex\ hsc)}
\leanline{4}{HasLaw\ (lpPhaseCoe\ (π.prod\ (kipldBrownianSource\ N\ dθ\ dx))}
\leanline{6}{(kipldBrownianFlow\ hN\ U\ L\ γ\ hdiff\ hLip\ π\ hπ\ T\ hT\ T))\ (gibbsPhaseTarget\ U)}
\leanline{6}{(π.prod\ (kipldBrownianSource\ N\ dθ\ dx))}
\end{leancode}
\leanlabel{How the Lean proof follows the argument}
\texttt{kipld\allowbreak{}Brownian\allowbreak{}Flow\_\allowbreak{}stationary} gives the formal counterpart of this step. The main intermediate facts are \texttt{\ensuremath{\pi}}, \texttt{h\ensuremath{\pi}}, \texttt{Q}, \texttt{Y}, \texttt{hl}, \texttt{hmarg} in that order. Named proof applications include \texttt{hp.\allowbreak{}symm}, \texttt{Prod.\allowbreak{}ext}; the adjacent proof shows their hypotheses and order of use.
\end{leanbox}


Earlier results used here: \cite[Section~3.26.13]{LangevinLean2026ULASC}; \cite[Section~3.26.15]{LangevinLean2026ULASC}; \cite[Section~3.26.3]{LangevinLean2026ULASC}; \cite[Section~3.1.1]{LangevinLean2026IPLA}; \cite[Section~3.1.3]{LangevinLean2026IPLA}; \cite[Section~3.1.4]{LangevinLean2026IPLA}; \cite[Section~3.5.1]{LangevinLean2026IPLA}; \cite[Section~3.17.1]{LangevinLean2026Kinetic}; \cite[Section~3.19.1]{LangevinLean2026Kinetic}; \cite[Section~3.33.1]{LangevinLean2026Kinetic}; \cite[Section~3.51.2]{LangevinLean2026Kinetic}; \cite[Section~3.58.3]{LangevinLean2026Kinetic}.

\begin{thebibliography}{ULA-LSI}
\bibitem[ULA-SC]{LangevinLean2026ULASC}
LangevinLean Project. \emph{Unadjusted Langevin under Strong Convexity}. LL-01, 2026. Draft report, content revision d34ead06987c. \url{https://odakyildiz.github.io/LangevinLean/docs/pdf/UnadjustedLangevin/FullProof.pdf}.
\bibitem[IPLA-SC]{LangevinLean2026IPLA}
LangevinLean Project. \emph{Interacting Particle Langevin under Strong Convexity}. LL-03, 2026. Draft report, content revision ec0201df2207. \url{https://odakyildiz.github.io/LangevinLean/docs/pdf/IPLA/FullProof.pdf}.
\bibitem[Kinetic]{LangevinLean2026Kinetic}
LangevinLean Project. \emph{Kinetic Langevin Monte Carlo}. LL-06, 2026. Draft report, content revision b2fee42a9593. \url{https://odakyildiz.github.io/LangevinLean/docs/pdf/KineticLangevin/FullProof.pdf}.
\end{thebibliography}

\end{document}
